%% kapa: exam=1 round=30 year=2019 period=1 form=A subject=civil_law subject_ko=민법 seq=8 number=8
%% source: Q-Net 감정평가사 기출문제 (cst00309) artlSeq=5205205 file=extracted/1008/5205205/2019년 제30회 감정평가사 1차시험 문제지/2019년 제30회 감정평가사 1차시험 문제지/2019년도 제30회 감정평가사 1차 1교시 A형.hwp
%% conversion: kordoc 4.24.51 -> tools/gen1.py (auto); review=proofread-sonnet; flags=none
\begin{question}{8}
법인 아닌 사단에 관한 설명으로 옳지 않은 것은? (다툼이 있으면 판례에 따름)
\begin{choices}
  \item 법인 아닌 사단의 사원이 집합체로서 물건을 소유할 때에는 총유로 한다.
  \item 법인 아닌 사단이 타인간의 금전채무를 보증하는 행위는 총유물의 관리 및 처분행위라고 볼 수 없다.
  \item 법인 아닌 사단의 총회 결의에 대해서는 민법상 사단법인에 대한 규정이 유추적용될 수 있다.
  \item 정관이나 규약에 정함이 없는 이상 사원총회의 결의를 거치지 않은 총유물의 관리 및 처분행위는 무효이다.
  \item 법인 아닌 사단은 부동산 등기능력이 없다.
\end{choices}
\end{question}
\answer{5}
